\documentclass[../report.tex]{subfiles}

\begin{document}
    \begin{abstract}
        This report presents a complete Interactive Reinforcement Learning (IRL) system for training a Pepper robot to perform contextually appropriate social greeting behaviours (DoNothing, Talk, Look, Wave, HandShake) in response to human cues. Three conditions are designed and evaluated: (1) an autonomous RL baseline using Proximal Policy Optimisation (PPO), (2) COACH fine-tuning with keyboard feedback, and (3) COACH fine-tuning with multimodal VR feedback. Conditions (2) and (3) independently fine-tune from the same pre-trained PPO checkpoint, allowing direct comparison of feedback modalities against each other and against the autonomous baseline. The PPO baseline uses a 12-dimensional observation space incorporating body-signal features (wrist height, wrist-to-core distance, body orientation, gaze direction), achieving a final training accuracy of 66.4\% and a test accuracy of 74.14\% after $\sim$150k training steps; a hyperparameter sweep across 10 configurations (30 runs total) identified learning rate ($\alpha = 0.001$) as the dominant factor. Because the limited human interaction budget is insufficient for PPO's batch-gradient requirements, both fine-tuning conditions use COACH (COrrective Advice Communicated by Humans), an algorithm that performs an immediate policy gradient update after every single feedback event. This targeted fine-tuning is designed to correct the complete Talk action suppression (0\% action distribution) introduced by feature overlap with the Look condition. The keyboard COACH session ran for 164 feedback steps and produced a measurable policy shift: action distribution entropy rose from 1.39 to 1.79 bits and HandShake occurrences increased 6.85$\times$. The VR condition---including controller buttons, head gestures (nod/shake), and voice commands fused via priority-based signal fusion---was fully implemented but could not be executed due to a physical body incompatibility between the NPC rig and the VR person controller during live sessions. Body-signal observations are normalised relative to each person's standing height, ensuring scale-invariant feature representations that generalise across the NPC rig and real human participants of differing heights. Evaluation compares all conditions on final task accuracy, action distribution entropy, and HandShake engagement as a proxy for policy diversity improvement.
    \end{abstract}
\end{document}